% howto/howto.tex
% mainfile: ../perfbook.tex
% SPDX-License-Identifier: CC-BY-SA-3.0

\QuickQuizChapter{chp:How To Use This Book}{如何使用本书}{qqzhowto}
%
\Epigraph{如果你能认识到生活是艰难的，事情对你来说反而会容易得多。}
	 {Louis D. Brandeis}

本书的目的是，在不损伤你神智的前提下，帮助你进行
shared-memory 并行系统编程。\footnote{
	或者更准确地说，所承受的理智风险不会比非并行编程大太多。
	不过仔细想想，这句话本身可能也说明不了什么。}
尽管如此，你应当把本书视为建筑的地基，
而非一座已经完工的大教堂。
你的使命——如果你选择接受的话——是帮助促进并行编程
这一激动人心的领域的发展，这些进展终有一天会使本书变得过时。

简言之，并行编程曾经被用于科学研究和具有超高挑战性的工程项目，
而现在它已迅速演化为一门工程学科。
因此，本书考察了一些与并行编程相关的具体问题，
并描述了解决这些问题的方法。
在某些出人意料的常见情况下，这些任务甚至可以被自动化。

编写本书时，笔者希望通过展示许多成功的并行编程项目背后的
工程方法论，让新一代高手从缓慢而艰辛地重新发明旧轮子的过程中
解放出来，将精力和创造力集中于新的前沿领域。
然而，你从本书中获得多少，取决于你投入多少。
我们希望仅仅阅读本书就会有所帮助，
而完成快速测验（Quick Quiz）则会更有帮助。
不过，最好的效果来自于将本书所教授的技术应用于实际问题。
一如既往，熟能生巧。

但无论你如何对待本书，我们都真切希望并行编程能给你带来
至少与给我们带来的同等多的乐趣、兴奋和挑战！

\section{路线图}
\label{sec:howto:Roadmap}
%
\epigraph{猫：你要去哪里？ \\
	  Alice：我该走哪条路？ \\
	  猫：那取决于你要去哪里。 \\
	  Alice：我不知道。 \\
	  猫：那走哪条路都无所谓了。}
	 {Lewis Carroll, \emph{Alice in Wonderland}}

本书是一本讲述通用且常用的设计技术的工具书，
而非一组只能解决特定问题的最优算法大全。
你目前正在阅读\cref{chp:How To Use This Book}，
不过这一点你已经知道了。
\Cref{chp:Introduction}给出了并行编程的高层次概述。

\Cref{chp:Hardware and its Habits}介绍了 shared-memory 并行硬件。
归根到底，不了解底层硬件就不可能写出好的并行代码。
由于硬件发展的速度很快，这一章难免有些过时。
不过我们会尽最大努力保证更新内容。
\Cref{chp:Tools of the Trade}随后提供了常用 shared-memory 并行编程原语的
简要概述。

\Cref{chp:Counting}深入研究了一个极其简单的问题的并行化——计数。
相信绝大多数读者都理解计数是怎么回事，所以本章能够在深入
各种重要的并行编程问题的同时，避免读者被问题本身所干扰。
据我了解，这一章在并行编程课程教学中使用最为广泛。

\Cref{chp:Partitioning and Synchronization Design}
引入了若干设计层面的方法，以解决\cref{chp:Counting}中
遇到的问题。
实践证明，在可行的情况下，在设计层面处理并行性可让人事半功倍：
借用 \pplsur{Edsger W.}{Dijkstra}~\cite{Dijkstra:1968:LEG:362929.362947}
的话来说，``后期加装的并行性被认为是极其次优的''~\cite{PaulEMcKenney2012HOTPARsuboptimal}。

接下来的三章考察了三种重要的同步手段。
\Cref{chp:Locking}涵盖了 locking，它不仅仍然是
生产级并行编程的主力，同时也被广泛认为是并行编程中最大的恶棍。
\Cref{chp:Data Ownership}简要概述了数据所有权（data ownership），
这是一种常被忽视但极为常用且有效的方法。
最后，\cref{chp:Deferred Processing}介绍了若干延迟处理机制，
包括 reference counting、hazard pointers、sequence locking 和 RCU\@。

\Cref{chp:Data Structures}将前几章学到的内容应用于 hash table。
Hash table 因其出色的可分区性而被大量使用，
其性能和可扩展性（通常）非常优秀。

不经验证的并行编程注定走向失败，相信很多读者对这一点
都刻骨铭心。
\Cref{chp:Validation}涵盖了各种形式的测试。
当然，光靠测试并不能证明你程序的可靠性，
因此\cref{chp:Formal Verification}紧接着简要概述了
几种实用的形式化验证方法。

\Cref{chp:Putting It All Together}
包含一系列中等规模的并行编程问题。
这些问题的难度各不相同，但对于已经熟练掌握前几章内容的读者来说
应该不成问题。

\Cref{sec:advsync:Advanced Synchronization}
探讨了高级同步方法，包括
non-blocking 同步和并行实时计算，
而\cref{chp:Advanced Synchronization: Memory Ordering}
涵盖了 memory ordering 这一高级主题。
\Cref{chp:Ease of Use}随后给出了一些易用性方面的建议。
\Cref{chp:Conflicting Visions of the Future}
展望了若干可能的未来方向，包括 shared-memory 并行系统设计、
软件和硬件 transactional memory，以及用于并行化的函数式编程。
最后，\cref{chp:Looking Forward and Back}回顾了本书的内容及其起源。

本书还有若干附录。
其中最受欢迎的似乎是
\cref{chp:app:whymb:Why Memory Barriers?}，
它更深入地探讨了 memory ordering。
\Cref{chp:app:Answers to Quick Quizzes}
包含了那些臭名昭著的快速测验的答案，我们将在下一节讨论这些测验。

\section{快速测验}
\label{sec:howto:Quick Quizzes}
%
\epigraph{去做困难的事情，否则你永远不会成长。}
	 {缩写自 Ronald E.~Osburn}

``快速测验''贯穿全书，答案可以在
\cref{chp:app:Answers to Quick Quizzes}中找到，
从第~\cpageref{chp:app:Answers to Quick Quizzes}~页开始。
其中一些基于该快速测验所在章节的内容，
但另一些则需要你思考超越该章节的范围，
在某些情况下，甚至超越当前知识的边界。
对于大多数读者来说，能从本书得到多少，
完全取决于为此付出的努力。
因此，在查看答案之前凭借自身努力解答测验的读者，
会发现他们的努力获得了丰厚的回报——对并行编程更深入的理解。

\QuickQuizSeries{%
\QuickQuizB{
	快速测验的答案在哪里可以找到？
}\QuickQuizAnswerB{
	在\cref{chp:app:Answers to Quick Quizzes}中，
	从第~\cpageref{chp:app:Answers to Quick Quizzes}~页开始。
	嗨，我觉得我应该给你出一道简单的！
}\QuickQuizEndB
%
\QuickQuizM{
	有些快速测验的问题似乎是从读者而非作者的视角提出的。
	这真的是有意为之的吗？
}\QuickQuizAnswerM{
	确实如此！
	许多问题是 Paul E.~McKenney 作为一个学习这些内容的初学者
	可能会问的问题。
	值得一提的是，Paul 的大部分并行编程知识是从并行硬件和软件中
	学到的，而非从教授那里学到的。
	根据 Paul 的经验，教授比并行系统更有可能回答口头提问，
	尽管近来语音助手取得了不少进步。
	当然，我们可以就教授和并行系统哪个能提供更有用的答案
	展开一场冗长的辩论，
	但就目前而言，让我们暂且同意，无论是教授还是并行系统，
	答案的有用程度在其群体中都差异巨大。

	其他一些测验与在涵盖本书内容的会议演讲和讲座中
	实际被问到的问题非常相似。
	还有一些则是从作者的视角出发的。
}\QuickQuizEndM
%
\QuickQuizE{
	这些快速测验实在不合我的口味。我能怎么办？
}\QuickQuizAnswerE{
以下是几种可能的策略：

\begin{enumerate}
\item	直接忽略快速测验，阅读本书的其余部分。
	你可能会错过某些快速测验中的有趣内容，
	但本书的其余部分同样有大量优秀内容。
	如果你的主要目标是获得对内容的总体了解，
	或者你正在浏览本书以寻找特定问题的解决方案，
	这是一种完全合理的方法。
\item	立即查看答案，而不是花大量时间自己想出答案。
	当某个快速测验的答案恰好是你正在尝试解决的
	特定问题的关键时，这种方法是合理的。
	如果你想对内容有更深入的理解，但不期望
	在仅凭一张白纸的情况下被要求生成并行解决方案，
	这种方法同样合理。
\item	如果你觉得快速测验令人分心但又无法忽略，
	你可以随时从 git 仓库克隆本书的 \LaTeX{} 源代码。\footnote{
		有关操作说明，请参见\cref{sec:howto:Whose Book Is This?}。}
	然后你可以运行命令 \co{make nq}，它会生成
	\co{perfbook-nq.pdf}。
	这个 PDF 在原本放置快速测验的位置会显示不显眼的方框标记，
	并将每章的快速测验收集到该章末尾，采用经典教科书的样式。
\item	学着喜欢（或至少容忍）快速测验。
	经验表明，在阅读时定期自我测验能极大地提高
	理解力和理解深度。
\end{enumerate}

请注意，快速测验与答案之间建立了超链接，反之亦然。
点击``快速测验''标题或小黑方块即可跳转到答案的开头。
从答案处，点击标题或小黑方块可跳转到测验的开头，
或者，点击答案末尾的小白方块可跳转到相应测验的末尾。
}\QuickQuizEndE
}

简言之，如果你想对本书的内容有深入的理解，就应该
花些时间在解答快速测验上。
别误会我的意思，被动阅读本书也是很有价值的，
但想拥有完整的问题解决能力，你就必须去实际解决问题。
同样，要获得完整的代码编写能力，确实需要你练习编写代码。

\QuickQuiz{
	如果被动阅读本书不能让我获得完整的问题解决和代码编写能力，
	那这本书究竟有什么意义？？？
}\QuickQuizAnswer{
	对于喜欢类比的人来说，编写并发软件就像演奏音乐一样，
	不同水平的才能和技能都有其用武之地。
	并非每个人都需要将毕生精力投入到成为一名
	音乐会钢琴家上。
	事实上，对于每一位这样的大师，都有大量水平较低的钢琴家，
	他们的音乐受到亲朋好友的欢迎。
	但这些水平较低的钢琴家可能在用其他方式谋生，
	并发编程也是如此。

	被动阅读本书的一个潜在好处是能够阅读和理解
	现代并发代码。这种能力反过来可以使你能够：

	\begin{enumerate}
	\item	了解并发代码的功能，从而检查某个提议的
		用例是否有效。
	\item	追踪并发代码中的 bug。
	\item	利用并发代码库中的信息，更容易地追踪与该代码库
		交互的其他代码中的 bug。
	\item	为并发代码中的 bug 编写修复补丁。
	\item	在并发代码库中创建简单直接的功能，无论是从头手写代码，
		还是使用现代的从互联网下载代码的
		``copy-pasta''开发方法论。\footnote{
			但请注意版权和许可问题！}
	\end{enumerate}

	如果你熟练掌握了 lock 和 atomic 操作的基本用法，
	被动阅读本书应该能使你成功地应用现代并发技术。

	最后，如果你的工作是协调使用现代并发技术的开发人员的活动，
	被动阅读本书可能有助于你理解他们到底在说什么。
}\QuickQuizEnd

我是在中年攻读博士学位期间的课程学习中才艰难地认识到这一点的。
当时我在学习一个熟悉的课题，
却惊讶地发现自己能脱口而出回答的章节练习题屈指可数。\footnote{
	所以我想，我的教授们拒绝让我免修那门课也许是件好事！}
强迫自己回答问题极大地加深了我对材料的记忆。
所以你看到的这些快速测验，我自己都先做了一遍！

最后，最常见的学习障碍是认为你已经理解了手头的材料。
快速测验可以成为一种极其有效的治疗方法。

\section{本书的替代读物}
\label{sec:Alternatives to This Book}
%
\epigraph{在两种邪恶之间，我总是选择我从未尝试过的那种。}
	 {Mae West}

正如 \pplsur{Donald}{Knuth} 在写书时所领悟到的那样，如果不想把一本书写得
无限长，那么就必须专注于某一方面。
本书专注于 shared-memory 并行编程，重点放在软件栈底层的软件，
例如操作系统内核、并行数据管理系统、底层库等。
本书使用的编程语言是 C。

如果你对并行性的其他方面感兴趣，可以参考其他书籍。
幸运的是，这样的书籍还真不少：

\begin{enumerate}
\item	如果你偏好更加学术、严格的并行编程处理方式，
	可能会喜欢 \pplsur{Maurice P.}{Herlihy} 和 \pplsur{Nir}{Shavit}
	的教科书~\cite{HerlihyShavit2008Textbook,HerlihyShavit2020Textbook}。
	这本书开头是一个有趣的搭配——底层原语与高层抽象的组合，
	之后是 locking 和简单数据结构（包括链表、队列、
	hash table 和计数器），最终以 transactional memory 收尾，
	全部使用 Java 语言。
	\ppl{Michael}{Scott} 的教科书~\cite{MichaelScott2013Textbook}
	以更偏向软件工程的视角处理类似的内容，
	据我所知，它是第一本正式出版的用整个章节讲述 RCU 的
	学术教科书。

	\pplsur{Maurice P.}{Herlihy}、\pplsur{Nir}{Shavit}、
	\pplsur{Victor}{Luchangco} 和 \pplsur{Michael}{Spear} 在
	第二版~\cite{HerlihyShavit2020Textbook}中赶了上来，
	增加了关于 hazard pointers 和 \acr{rcu} 的简短章节，
	后者以 \acr{ebr} 的名义出现。\footnote{
		尽管其中包含一个由 \ppl{Richard}{Bornat} 指出的
		reader-preemption bug。}
	他们还包含了两者的简史，不过 \acr{rcu} 的历史
	从其被 Linux 内核接受近一年后才开始叙述，
	距离 \pplsur{H. T.}{Kung} 和 \pplsur{Philip L.}{Lehman}
	的里程碑论文~\cite{Kung80}已有 20 多年。
	希望深入了解历史的读者可以参阅
	本书的\cref{sec:defer:RCU Related Work}。

	然而，那些可能怀疑该教科书第一作者对 RCU 持敌意态度的读者，
	应参考其一篇论文~\cite{Balmau:2016:FRM:2935764.2935790}
	第一页最后一个完整句子。
	这个句子写道：``QSBR [一类特定的 RCU 实现] 速度快，
	可以应用于几乎任何数据结构。''
	这显然不是一个对 RCU 持敌意态度的人会说的话。
\item	如果你喜欢从偏学术的编程语言学角度了解并行编程，
	那么可能会对 \pplsur{Michael}{Scott} 的编程语言实用主义
	教科书~\cite{MichaelScott2006Textbook,MichaelScott2015Textbook}
	中的并发章节感兴趣。
\item	如果你喜欢用面向对象的 C++ 来了解并行编程，
	可以试试 \pplsur{Douglas C.}{Schmidt} 的 POSA 系列的
	第~2 和第~4 卷~\cite{SchmidtStalRohnertBuschmann2000v2Textbook,
	BuschmannHenneySchmidt2007v4Textbook}。
	第~4 卷有几章有意思的内容，用 C++ 编写了一个仓库管理软件。
	其中名为``分割大泥巴球''的章节标题就可以看出该示例的真实性，
	并行性本身的问题反而退居次要地位，让位于理解现实世界应用的复杂性。
\item	如果你想把并行编程应用于 Linux 内核设备驱动程序，
	那么 \pplsur{Jonathan}{Corbet}、\pplsur{Alessandro}{Rubini}
	和 \pplsur{Greg}{Kroah-Hartman} 的
	``Linux Device Drivers''~\cite{CorbetRubiniKroahHartman}
	绝对是必选项，Linux Weekly News 网站
	(\url{https://lwn.net/}) 同样如此。
	关于 Linux 内核内部机制这一更广泛的主题，
	有大量的书籍和网上资源可供参考。
\item	如果你的主要目标是科学和技术计算，
	同时偏好模式化方法，
	可以看看 \pplsur{Timothy G.}{Mattson} 等人的
	教科书~\cite{Mattson2005Textbook}。
	它涵盖了 Java、C/C++、OpenMP 和 MPI\@。
	这本书采用的理念是设计优先，实现第二。
\item	如果你的主要目标是科学和技术计算，
	同时对 GPU、CUDA 和 MPI 感兴趣，
	可以查阅 \ppl{Norm}{Matloff} 的
	``Programming on Parallel Machines''~\cite{NormMatloff2017ParProcBook}。
	当然，GPU 厂商也有相当多的额外
	资料~\cite{AMD2020ROCm,CyrilZeller2011GPGPUbasics,NVidia2017GPGPU,NVidia2017GPGPU-university}。
\item	如果你主要关注低延迟，尤其是使用 Rust 语言的话，
	\ppl{Pekka}{Enberg} 的
	``Latency: Reduce delay in software
	systems''~\cite{PekkaEnberg2026Latency}
	可能正是你需要的。
\item	如果你对 POSIX Threads 感兴趣，
	可以看看 \pplmdl{David R.}{Butenhof} 的书~\cite{Butenhof1997pthreads}。
	此外，\ppl{W.~Richard}{Stevens} 的
	不朽之作~\cite{WRichardStevens1992,WRichardStevens2013}
	涵盖了 UNIX 和 POSIX，\ppl{Stewart}{Weiss} 的讲义
	~\cite{StewartWeiss2013UNIX}详尽又通俗地介绍了很多
	相关的示例。
\item	如果你对 C++11 感兴趣，你可能会喜欢
	\ppl{Anthony}{Williams} 的``C++ Concurrency in Action:
	Practical Multithreading''~\cite{AnthonyWilliams2012,AnthonyWilliams2019}
	或 \ppl{Rainer}{Grimm} 的``Concurrency with Modern
	C++''~\cite{RainerGrimm2017C++Concurrency}。
\item	如果你对 C++ 感兴趣，但是在 Windows 环境中，
	可以尝试 \ppl{Herb}{Sutter} 的``Effective Concurrency''
	系列，发表在
	Dr.~Dobbs Journal~\cite{HerbSutter2008EffectiveConcurrency}上。
	该系列在展示常识性的并行编程方法方面做得相当不错。
\item	如果你想尝试 Intel Threading Building Blocks，
	那么 \ppl{James}{Reinders} 的书~\cite{Reinders2007Textbook}
	可能正是你要找的。
\item	对于有兴趣了解各种类型的多处理器硬件缓存组织
	如何影响内核内部实现的读者，
	应该看看 \ppl{Curt}{Schimmel} 对这一主题的
	经典论述~\cite{Schimmel:1994:USM:175689}。
\item	如果你寻找硬件视角，\pplsur{John L.}{Hennessy} 和
	\pplsur{David A.}{Patterson} 的经典
	教科书~\cite{Hennessy2017}非常值得一读。
	面向科学和技术工作负载（处理大数组）的精简版可以在
	\ppl{Andrew}{Chien} 的
	教科书~\cite{AndrewChien2022ComputerArchitectureScientists}中找到。
	如果你寻找从更偏硬件视角论述 memory ordering 的学术教科书，
	\ppl{Daniel}{Sorin} 等人的~\cite{DanielJSorin2011MemModel,%
	VijayNagarajan2020MemModel}
	强烈推荐。
	关于从 Linux 内核视角介绍 memory ordering 的教程，
	\ppl{Paolo}{Bonzini} 的 LWN 系列是一个好的起点~\cite{PaoloBonzini2021lockless1,PaoloBonzini2021lockless2,PaoloBonzini2021lockless3,PaoloBonzini2021lockless4,PaoloBonzini2021lockless5,PaoloBonzini2021lockless6}。
\item	希望了解 Rust 语言对底层并发支持的读者应参考
	\ppl{Mara}{Bos} 的书~\cite{MaraBos2023RustAtomicsAndLocks}。
\item	最后，使用 Java 的读者可能会从 \ppl{Doug}{Lea} 的
	教科书~\cite{DougLea1997Textbook,Goetz2007Textbook}中受益匪浅。
\end{enumerate}

然而，如果你想阅读一本讲述底层软件（尤其是用 C 编写的软件）的
并行设计原则的书，请继续阅读！

\section{示例源代码}
\label{sec:howto:Sample Source Code}
%
\epigraph{用源代码，Luke！}{不知名的 Star Wars 粉丝}

本书有相当多的内容在讨论源代码，多数时候书中的源代码可以在
本书 git 仓库的 \path{CodeSamples} 目录中找到。
例如，在 UNIX 系统上，你可以输入以下命令：

\begin{VerbatimU}
find CodeSamples -name rcu_rcpls.c -print
\end{VerbatimU}

这个命令将找到文件 \path{rcu_rcpls.c}，它在
\cref{chp:app:``Toy'' RCU Implementations}中被引用。
非 UNIX 系统有其自己众所周知的按文件名查找文件的方法。

\section{视频资源}
\label{sec:howto:Video Resources}
%
\epigraph{我听了就忘了；我看了就记住了；我做了就理解了。}
	 {佚名}
% Attributed to Confucius or Xunzi.  Or to the 1960s USA experiential
% education movement.  Who knows?

本节介绍了一些涵盖本书部分章节的演讲。

\Cref{chp:Introduction}（``Introduction''）由 2025 年题为
``Breaking Up is Hard to Do''的演讲涵盖。
视频可在此处获取：
\url{https://kernel-recipes.org/en/2025/schedule/breaking-up-is-hard-to-do/}，
但你需要自备 2x4 的 Duplo 和 Lego 积木，前者两块，后者十六块。
该演讲还涉及了
\cref{chp:Partitioning and Synchronization Design}
（``Partitioning and Synchronization Design''）。

\Cref{chp:Hardware and its Habits}（``Hardware and its Habits''）
由 2023 年题为``Hardware and its Concurrency Habits''的演讲涵盖。
幻灯片和视频可在此处获取：
\url{https://kernel-recipes.org/en/2023/schedule/hardware-and-its-concurrency-habits/}。

\Cref{chp:Counting}（``Counting''）由 2024 年题为``Case Study:
Concurrent Counting''的演讲涵盖。
幻灯片和视频可在此处获取：
\url{https://kernel-recipes.org/en/2024/schedule/case-study-concurrent-counting/}。

\Cref{sec:defer:Hazard Pointers}（``Hazard Pointers''）由
2024 年题为``Hazard Pointers in the Linux Kernel''的演讲涵盖，
尽管我们希望 Maged Michael 有朝一日能做一次有视频存档的演讲。
与此同时，幻灯片和视频可在此处获取：
\url{https://lpc.events/event/18/contributions/1731/}。

\Cref{sec:defer:Introduction to RCU}（``Introduction to RCU''）
和
\cref{sec:defer:RCU Fundamentals}（``RCU Fundamentals''）
由 2021 年题为``Unraveling RCU-Usage Mysteries
(Fundamentals)''的演讲涵盖。
幻灯片和视频分别可在此处获取：
\url{http://www.rdrop.com/~paulmck/RCU/RCUusageFundamental.2021.12.07a.LF.pdf}
和
\url{https://www.linuxfoundation.org/webinars/unraveling-rcu-usage-mysteries}。

\Cref{sec:defer:RCU Usage}（``RCU Usage''）
由 2022 年题为``Unraveling RCU-Usage Mysteries
(Additional Use Cases)''的演讲涵盖。
幻灯片和视频分别可在此处获取：
\url{https://events.linuxfoundation.org/wp-content/uploads/2022/02/RCUusageAdditional.2022.02.22b.LF-1.pdf}
和
\url{https://www.linuxfoundation.org/webinars/unraveling-rcu-usage-mysteries-additional-use-cases}。

\Cref{sec:debugging:Hunting Heisenbugs}（``Hunting Heisenbugs''）
由 2023 年同样题为``Hunting Heisenbugs''的演讲涵盖。
幻灯片和视频可在此处获取：
\url{https://lpc.events/event/17/contributions/1504/}。

\section{这是谁的书？}
\label{sec:howto:Whose Book Is This?}
%
\epigraph{如果你成为一名教师，你的学生们会教导你。}
	 {Oscar Hammerstein II}

正如封面所写，编著者是 Paul E.~McKenney。
不过，编著者确实通过
\href{mailto:perfbook@vger.kernel.org}
{\nolinkurl{perfbook@vger.kernel.org}} 邮件列表接受投稿。
投稿可以是任何形式，常见的比如文本邮件、
针对本书 \LaTeX{} 源代码的补丁，或者 \co{git pull} 请求。
请使用任何对你来说方便的方式。

想要创建补丁或 \co{git pull} 请求，你需要本书的
\LaTeX{} 源代码，位于
\url{git://git.kernel.org/pub/scm/linux/kernel/git/paulmck/perfbook.git}，
或者
\url{https://git.kernel.org/pub/scm/linux/kernel/git/paulmck/perfbook.git}。
此外你还需要 \co{git} 和 \LaTeX{}，在大多数主流 Linux 发行版中
都可以找到它们。
在某些发行版里，你还需要其他的安装包。
几个常用发行版所需的安装包列表列在
本书 \LaTeX{} 源代码中的 \path{FAQ-BUILD.txt} 文件中。

\begin{listing}
\begin{VerbatimL}[breaklines=true,breakafter=/,
        breakaftersymbolpre=\raisebox{-.7ex}{\textcolor{darkgray}{\Pisymbol{psy}{191}}},
	breaksymbolleft=\textcolor{darkgray}{\tiny\ensuremath{\hookrightarrow}},
        numbers=none,xleftmargin=0pt]
git clone git://git.kernel.org/pub/scm/linux/kernel/git/paulmck/perfbook.git
cd perfbook
# You may need to install a font. See item 1 in FAQ.txt.
make                     # -jN for parallel build
evince perfbook.pdf &    # Two-column version
make perfbook-1c.pdf
evince perfbook-1c.pdf & # One-column version for e-readers
make help                # Display other build options
\end{VerbatimL}
\caption{创建最新的 PDF}
\label{lst:howto:Creating a Up-To-Date PDF}
\end{listing}

想要创建并显示本书当前的 \LaTeX{} 源代码，
请使用\cref{lst:howto:Creating a Up-To-Date PDF}所示的
Linux 命令。
在某些环境中，显示 \path{perfbook.pdf} 的 \co{evince} 命令
可能需要替换为其他命令，例如 \co{acroread}。
\co{git clone} 命令只在第一次生成 PDF 时需要，
之后只需运行\cref{lst:howto:Generating an Updated PDF}
中的命令来更新代码并生成更新后的 PDF\@。
\cref{lst:howto:Generating an Updated PDF}中的命令
必须在\cref{lst:howto:Creating a Up-To-Date PDF}所示命令
创建的 \path{perfbook} 目录内运行。

\begin{listing}
\begin{VerbatimL}[numbers=none,xleftmargin=0pt]
git remote update
git checkout origin/master
make                     # -jN for parallel build
evince perfbook.pdf &    # Two-column version
make perfbook-1c.pdf
evince perfbook-1c.pdf & # One-column version for e-readers
\end{VerbatimL}
\caption{生成更新后的 PDF}
\label{lst:howto:Generating an Updated PDF}
\end{listing}

本书的 PDF 不时会发布在
\url{https://kernel.org/pub/linux/kernel/people/paulmck/perfbook/perfbook.html}
和
\url{http://www.rdrop.com/users/paulmck/perfbook/}。

给本书提交补丁和发送 \co{git pull} 请求的流程与 Linux 内核非常类似，
其文档在此处：
\url{https://www.kernel.org/doc/html/latest/process/submitting-patches.html}。
其中一个重要要求是，每个补丁（或 \co{git pull} 请求中的每个提交）
必须包含一条有效的 \co{Signed-off-by:} 行，格式如下：

\begin{VerbatimU}
Signed-off-by: My Name <myname@example.org>
\end{VerbatimU}

请参阅 \url{https://lore.kernel.org/lkml/20070116022324.GA28513@linux.vnet.ibm.com/}
获取一个包含 \co{Signed-off-by:} 行的补丁示例。
请注意，\co{Signed-off-by:} 行有着非常明确的含义，
即你在保证：

\begin{enumerate}[label={(\alph*)}]
\item	贡献内容全部或部分由我创建，我有权根据文件中
	指明的开源许可证提交它；或
\item	贡献内容基于先前的工作，据我所知，这些工作受
	适当的开源许可证保护，我有权将由我全部或部分创建的
	修改后的工作以相同的开源许可证提交
	（除非我被允许使用不同的许可证），
	如文件中所示；或
\item	贡献由其他人直接提供给我，该人已保证 (a)、(b) 或 (c)，
	且我没有修改任何内容。
\item	我理解并同意该项目和贡献是公开的，
	贡献的记录（包括我随之提交的所有个人信息，
	包括我的签名）将被无限期保留，
	这些信息可以随着本项目或相关的开源许可证被再分发。
\end{enumerate}

这与 Linux 内核使用的 Developer's Certificate of Origin (DCO) 1.1
非常相似。
必须使用真实姓名：
遗憾的是，我无法接受匿名或化名的贡献。

本书的写作语言是美式英语，不过本书的开源性质允许翻译，
我个人也鼓励大家翻译本书。
覆盖本书的开源许可证还允许你出售你的译本（如果你愿意的话）。
我只希望你给我寄一份译本（如有精装本则更好），
但这只是出于职业礼貌的请求，
绝不是你在 Creative Commons 和 GPL 许可证下已经拥有的
权限的前提条件。
在源代码目录的 \co{FAQ.txt} 中列出了当前正在进行的翻译。
我认可的``进行中''是至少翻译了一章以上的内容。

我没有计划出版本书的精装本，但你可以随意这样做，
例如使用互联网上众多印刷服务中的任何一种。
封面艺术可在本书 \co{git} 仓库的 \path{cartoons} 目录中获得，
但你也可以自由设计自己的封面。

``美式英语''这一大类下有许多风格。
本书的具体风格记录在\cref{chp:app:styleguide:Style Guide}中。

正如我在一开始所说的，我只是本书的编著者。
如果你愿意对本书做出贡献，这也会是你的书。
本着这种精神，我将为你呈现\cref{chp:Introduction}，我们的引言。

\QuickQuizAnswersChp{qqzhowto}
